% ID: SEL-D0063-OA-E01
% ESTADO: cotejada. Segunda pasada visual y borrador compilado revisados.
% Fuente: Selectividad/sel_2010_matematicas_aplicadas/mat_apli/mat_apli_junio/mat_apli_junio_rb/EXAMEN-ANDALUCÍA-6.pdf
% SHA-256: 0de2f7d077a3b5d5c0dbca886cdb363082356411843369c3bce54ff7cda08abc
% Página: 1; opción A; ejercicio original 1.

\textbf{(2.5 puntos)} Un supermercado se abastece de gambas y langostinos a través de dos mayoristas, A y B, que le envían contenedores con cajas completas de ambos productos.

El mayorista A envía en cada contenedor 2 cajas de gambas y 3 de langostinos, al precio de 350 euros el contenedor, mientras que el mayorista B envía en cada uno 1 caja de gambas y 5 de langostinos, al precio de 550 euros el contenedor.

El supermercado necesita, como mínimo, 50 cajas de gambas y 180 de langostinos pudiendo almacenar, como máximo, 50 contenedores.

\textquestiondown{}Cuántos contenedores debería pedir el supermercado a cada mayorista para satisfacer sus necesidades con el menor coste posible? Indique cuál sería ese coste mínimo.
